\documentclass[12pt]{article}
\usepackage[a4paper, margin=2.5cm]{geometry} % Indstiller margen til 2.5 cm
\usepackage{setspace} % Til linjeafstand
\onehalfspacing % Indstiller linjeafstand til 1.5
\usepackage{graphicx} % Til billeder
\usepackage{amsmath, amssymb} % Til matematiske udtryk
\usepackage{booktabs} % Til tabeller
\usepackage{caption} % Bedre kontrol over billedtekst
\usepackage{float} % Tillader placering af figurer på en ny side
\usepackage{pgfplots}
\usetikzlibrary{intersections, patterns, decorations.markings, calc}

\title{Macroeconomics III - Assignment 1}
\author{fps618, kvb125}
\date{feb 2026}

\begin{document}

\maketitle

\section{Exercise 1}
\subsection{a)}
We want to find the optimal consumption path meaning how much to consume each period. To do this, we set up a Lagrangian and take derivatives. Our constraint each period is: $c_t+a_{t+1}=y_t+(1+r) a_t$, which we rearrange: $y_t+(1+r) a_t-c_t-a_{t+1}=0$ Since this constraint holds for every period, we need a separate multiplier $\lambda_t$ for each one: 
$$
\mathcal{L}=\sum_{t=0}^{\infty} \beta^t \frac{c_t^{1-\sigma}}{1-\sigma}_{\text {utility }}+\sum_{t=0}^{\infty} \lambda_t\left[y_t+(1+r) a_t-c_t-a_{t+1}\right]
$$
We have two choice variables each period: $\mathbf{c}_{\mathbf{t}}$ (consumption) and $\mathbf{a}_{\mathbf{t}_{+1}}$ (savings) so we find FOCs with respect for both. 
\subsubsection{FOC with respect to $c_t$}
$$
\frac{\partial \mathcal{L}}{\partial c_t}=\beta^t \cdot \frac{(1-\sigma) c_t^{-\sigma}}{1-\sigma}+\lambda_t \cdot(-1)=\beta^t c_t^{-\sigma}-\lambda_t
$$
Setting the sum of both terms equal to zero:
$$
\frac{\partial \mathcal{L}}{\partial c_t}=\beta^t c_t^{-\sigma}-\lambda_t=0
$$
$$
\boxed{\lambda_t=\beta^t c_t^{-\sigma}}
$$

\subsubsection{FOC with respect to $a_{t_{+1}}$}
Because $\mathbf{a}_{\mathbf{t}_{\mathbf{+ 1}}}$ appears in two different period constraints, we need to differentiate term by term:
\\
\textbf{Period $t$}:
$$
\frac{\partial}{\partial a_{t+1}} \lambda_t\left[y_t+(1+r) a_t-c_t-a_{t+1}\right]=\lambda_t \cdot(-1)=-\lambda_t
$$
\textbf{Period $t+1$}
$$
\frac{\partial}{\partial a_{t+1}} \lambda_{t+1}\left[y_{t+1}+(1+r) a_{t+1}-c_{t+1}-a_{t+2}\right]=\lambda_{t+1} \cdot(1+r)
$$
Setting the sum of both terms equal to zero gives us the second FOC:
$$
\frac{\partial \mathcal{L}}{\partial a_{t+1}}=-\lambda_t+(1+r) \lambda_{t+1}=0
$$
$$\boxed{\lambda_t=(1+r) \lambda_{t+1}}$$
\subsubsection{Obtaining the euler-equation}
To better understand the trade off between consumption and saving, we can substitute the $c_t$ FOC into the $a_t+t$ FOC. 
\\
\\
From the $c_t$ FOC we got:
$$
\lambda_t=\beta^t c_t^{-\sigma}
$$
This holds for any period as it's a general statement. So if we shift it forward one period we get:
$$
\lambda_{t+1}=\beta^{t+1} c_{t+1}^{-\sigma}
$$
These two terms for $\lambda_t$ and $\lambda_t+1$ can now be inserted directly into the second FOC:
$$
\underbrace{\beta^t c_t^{-\sigma}}_{\lambda_t}=(1+r) \cdot \underbrace{\beta^{t+1} c_{t+1}^{-\sigma}}_{\lambda_{t+1}}
$$
Finally we divide with $\beta^t$ on both sides to arrive at the euler equation:
$$
\boxed{c_t^{-\sigma}=\beta(1+r) \cdot c_{t+1}^{-\sigma}}
$$
The Euler equation tells us exactly what influences the tradeoff between consuming today vs. tomorrow adjusted for impatience $\beta$ and interest $1+r$. $\sigma$ is an elasticity parameter. A higher $\sigma$ means you're less willing to substitute across time even if $r$ is high, you won't shift much consumption to the future. 
\\
\\
At the optimum, you're indifferent between consuming one unit today or saving it and consuming (1+r) units tomorrow. If this weren't true, you could make yourself better off by shifting consumption between periods $\rightarrow$ thus you have not found the optimal bundle.

\subsection{b)}

\end{document}